\documentclass[12pt,a4paper]{article}
\usepackage[utf8]{inputenc}
\usepackage[T1]{fontenc}
\usepackage[spanish]{babel}
\usepackage{mathpazo} % Palatino para texto y matemáticas
\usepackage{amsmath, amssymb, amsthm}
\usepackage{geometry}
\geometry{margin=2.5cm}
\usepackage{parskip}

\title{\textbf{Cambio de variable en integrales de una dimensión}}
\author{Franz Chouly}
\date{\today}

\begin{document}

\maketitle

\section{Preliminares: integración con límites invertidos}

En esta sección se discute cómo el \textbf{Teorema Fundamental del Cálculo} permite dar sentido a integrales donde el límite inferior es mayor que el superior, es decir, integrales de la forma \(\int_b^a f(t) \, dt\) con \(b < a\).

\subsection{Teorema fundamental del cálculo}
El \textbf{Teorema Fundamental del Cálculo} establece que si \(f\) es una función \textbf{continua} en un intervalo \([a, b]\), entonces existe una función \(F\) (llamada \textbf{primitiva} o \textbf{antiderivada} de \(f\)) tal que:
\[
F'(x) = f(x) \quad \text{para todo } x \in [a, b],
\]
y además,
\[
\int_a^b f(t) \, dt = F(b) - F(a).
\]

\subsection{Inversión de los límites de integración}
Supongamos ahora que \(b < a\). Para definir \(\int_b^a f(t) \, dt\), utilizamos la primitiva \(F\) de \(f\) y aplicamos el teorema fundamental de manera consistente:
\[
\int_b^a f(t) \, dt = F(a) - F(b).
\]
Sin embargo, observamos que:
\[
F(a) - F(b) = - \left( F(b) - F(a) \right) = - \int_a^b f(t) \, dt.
\]
Por lo tanto, podemos \textbf{definir} la integral con límites invertidos como:
\[
\int_b^a f(t) \, dt = - \int_a^b f(t) \, dt.
\]

\subsection{Coherencia con la propiedad de aditividad (Regla de Chasles)}
Esta definición es \textbf{coherente} con la \textbf{propiedad de aditividad} de las integrales (también conocida como \textbf{regla de Chasles}), que establece que para cualquier \(a, b, c\) en el dominio de \(f\):
\[
\int_a^c f(t) \, dt = \int_a^b f(t) \, dt + \int_b^c f(t) \, dt.
\]

En particular, si tomamos \(c = a\), obtenemos:
\[
\int_a^a f(t) \, dt = \int_a^b f(t) \, dt + \int_b^a f(t) \, dt.
\]
Pero \(\int_a^a f(t) \, dt = 0\) (la integral sobre un intervalo de longitud cero es nula). Por lo tanto:
\[
0 = \int_a^b f(t) \, dt + \int_b^a f(t) \, dt.
\]
Despejando, recuperamos la definición anterior:
\[
\int_b^a f(t) \, dt = - \int_a^b f(t) \, dt.
\]

\subsection{Interpretación geométrica}
Desde un punto de vista geométrico, la integral \(\int_a^b f(t) \, dt\) representa el \textbf{área con signo} bajo la curva \(y = f(t)\) entre \(a\) y \(b\):
\begin{itemize}
    \item Si \(a < b\), el área se calcula de izquierda a derecha.
    \item Si \(b < a\), la integral \(\int_b^a f(t) \, dt\) invierte el sentido del recorrido (de derecha a izquierda), lo que justifica el signo negativo.
\end{itemize}

\subsection{Ejemplo}
Consideremos \(f(t) = t^2\) y calculemos \(\int_2^1 t^2 \, dt\):
\[
\int_2^1 t^2 \, dt = - \int_1^2 t^2 \, dt = - \left[ \frac{t^3}{3} \right]_1^2 = - \left( \frac{8}{3} - \frac{1}{3} \right) = - \frac{7}{3}.
\]
Esto coincide con la primitiva evaluada directamente:
\[
F(t) = \frac{t^3}{3}, \quad \int_2^1 t^2 \, dt = F(1) - F(2) = \frac{1}{3} - \frac{8}{3} = - \frac{7}{3}.
\]

\subsection{Observación final}
Esta convención es \textbf{universal} en matemáticas y física, y garantiza que las propiedades algebraicas de las integrales (como la linealidad y la aditividad) se mantengan incluso cuando los límites están invertidos.

%***
%\section{Preliminar}

%Teorema fundamental, si $b<a$, el teorema fundamental permite dar un sentido a
%\[
%\int_b^a f(t) dt.
%\]
%En efecto
%\[
%\int_b^a f(t) dt = F(a) - F(b) = - (F(b) - F(a))
%=:
%- \int_a^b f(t) dt.
%\]
%Es coherente con la regla de Chasles, porque tenemos
%\[
%0= \int_a^a f(t)dt 
%= \int_a^b f(t)dt + \int_b^a f(t)dt.
%\]

\section{Fórmula de cambio de variable en $\mathbb{R}^1$}


Dada una función $f: [a, b] \to \mathbb{R}$ integrable, y una función $\phi: [\alpha, \beta] \to [a, b]$ \textbf{biyectiva}, \textbf{continua} y \textbf{diferenciable} con $\phi'(\theta) \neq 0$ para todo $\theta \in [\alpha, \beta]$, se cumple:
\[
\int_{a}^{b} f(x) \, dx = \int_{\alpha}^{\beta} f(\phi(\theta)) \cdot |\phi'(\theta)| \, d\theta
\]

\subsection*{Explicación}
\begin{itemize}
    \item $\phi'(\theta)$ es la \textbf{derivada} de $\phi$ (el "jacobiano" en 1D).
    \item El valor absoluto $|\phi'(\theta)|$ asegura que la integral preserve el signo del área.
%    \item Si $\phi$ es \textbf{decreciente} (es decir, $\phi'(\theta) < 0$), los límites de integración se invierten: $\alpha \to \beta$ pasa a ser $\beta \to \alpha$.
\end{itemize}

\subsection{Demostración del teorema de cambio de variable}

\paragraph{Caso $\phi$ creciente}
Supongamos que $\phi: [\alpha, \beta] \to [a, b]$ es una función de clase $C^1$ (continuamente diferenciable) y \textbf{estrictamente creciente}, con $\phi'(\theta) > 0$ para todo $\theta \in [\alpha, \beta]$. Sea $f: [a, b] \to \mathbb{R}$ una función continua, y $F$ una primitiva de $f$ en $[a, b]$, es decir, $F'(x) = f(x)$ para todo $x \in [a, b]$.

Aplicamos el \textbf{Teorema Fundamental del Cálculo} y la \textbf{regla de la cadena}:
\begin{align*}
\int_{a}^{b} f(x) \, dx &= F(b) - F(a) \\
&= F(\phi(\beta)) - F(\phi(\alpha)) \quad \text{(por definición de $a$ y $b$)} \\
&= \int_{\alpha}^{\beta} \frac{d}{d\theta} [F(\phi(\theta))] \, d\theta \quad \text{(Teorema Fundamental)} \\
&= \int_{\alpha}^{\beta} F'(\phi(\theta)) \cdot \phi'(\theta) \, d\theta \quad \text{(Regla de la Cadena)} \\
&= \int_{\alpha}^{\beta} f(\phi(\theta)) \cdot \phi'(\theta) \, d\theta \quad \text{(pues $F' = f$)}.
\end{align*}

\paragraph{Caso $\phi$ decreciente}
Ahora, supongamos que $\phi: [\alpha, \beta] \to [a, b]$ es de clase $C^1$ y \textbf{estrictamente decreciente}, con $\phi'(\theta) < 0$ para todo $\theta \in [\alpha, \beta]$. En este caso, los extremos del intervalo se invierten: $a = \phi(\beta)$ y $b = \phi(\alpha)$.

Repitiendo el procedimiento:
\begin{align*}
\int_{a}^{b} f(x) \, dx &= F(b) - F(a) \\
&= F(\phi(\alpha)) - F(\phi(\beta)) \quad \text{(pues $a = \phi(\beta)$ y $b = \phi(\alpha)$)} \\
&= - \left( F(\phi(\beta)) - F(\phi(\alpha)) \right) \\
&= - \int_{\alpha}^{\beta} \frac{d}{d\theta} [F(\phi(\theta))] \, d\theta \\
&= - \int_{\alpha}^{\beta} f(\phi(\theta)) \cdot \phi'(\theta) \, d\theta \\
&= \int_{\beta}^{\alpha} f(\phi(\theta)) \cdot \phi'(\theta) \, d\theta \quad \text{(inversión de límites)} \\
&= \int_{\alpha}^{\beta} f(\phi(\theta)) \cdot |\phi'(\theta)| \, d\theta \quad \text{(pues $\phi'(\theta) < 0$)}.
\end{align*}

\subsection{Síntesis}
En ambos casos, el cambio de variable en integrales se puede expresar de manera unificada como:
\[
\int_{a}^{b} f(x) \, dx = \int_{\alpha}^{\beta} f(\phi(\theta)) \cdot |\phi'(\theta)| \, d\theta,
\]
donde $\phi$ es una función biyectiva y de clase $C^1$ en $[\alpha, \beta]$. El valor absoluto en $|\phi'(\theta)|$ garantiza que la integral conserve su signo independientemente de si $\phi$ es creciente o decreciente.

%***
%\subsection{Demontracion}

%\paragraph{Caso $\phi$ creciente}

%En este caso, tenemos
%\[
%a = \phi(\alpha)\quad b = \phi(\beta).
%\]
%Entonces, con $F$ primitiva de $f$ y aplicando el Teorema Fundamental y luego la regla de la cadena y luego el teorema fundamental de nuevo.
%\begin{eqnarray}
%\int_{a}^{b} f(x) \, dx 
%&=& F(b) - F(a) \\
%&=& F(\phi(\beta)) - F(\phi(\alpha))\\
%&=& \int_\alpha^\beta (F \circ \phi)'(\theta)
%d\theta\\
%&=& \int_\alpha^\beta (F' \circ \phi)(\theta) \phi'(\theta)
%d\theta\\
%&=& \int_\alpha^\beta f ( \phi(\theta)) \phi'(\theta)
%d\theta
%.
%\end{eqnarray}

%Caso D
%
%Resumimos
%$I = (a;b)$ $J=\min$
%entonces


%\subsection{Mas detalles}

\subsection{Ejemplo clásico: sustitución trigonométrica}
Para calcular $\int_{-1}^{1} \frac{1}{\sqrt{1 - x^2}} \, dx$, usamos $x = \sin(\theta)$:
\begin{align*}
    \phi(\theta) &= \sin(\theta), \quad \phi'(\theta) = \cos(\theta), \\
    \phi^{-1}(-1) &= -\frac{\pi}{2}, \quad \phi^{-1}(1) = \frac{\pi}{2}.
\end{align*}
La integral se transforma en:
\[
\int_{-\pi/2}^{\pi/2} \frac{1}{\sqrt{1 - \sin^2(\theta)}} \cdot \cos(\theta) \, d\theta = \int_{-\pi/2}^{\pi/2} 1 \, d\theta = \pi.
\]

\section{Aplicación a funciones pares e impares}
Si $f$ es una función \textbf{par} ($f(-x) = f(x)$) o \textbf{impar} ($f(-x) = -f(x)$), y el intervalo $[a, b]$ es simétrico respecto al origen (es decir, $a = -b$), entonces:

\[
\int_{-b}^{b} f(x) \, dx =
\begin{cases}
2 \displaystyle \int_{0}^{b} f(x) \, dx & \text{si } f \text{ es par}, \\
0 & \text{si } f \text{ es impar}.
\end{cases}
\]

\subsection*{Demostración para funciones impares}
Usamos el cambio de variable $x = -u$:
\begin{align*}
    \int_{-b}^{b} f(x) \, dx &= \int_{b}^{-b} f(-u) \cdot (-1) \, du \quad (\text{porque } dx = -du) \\
    &= -\int_{b}^{-b} f(-u) \, du = \int_{-b}^{b} f(-u) \, du.
\end{align*}
Si $f$ es impar, $f(-u) = -f(u)$, por lo que:
\[
\int_{-b}^{b} f(x) \, dx = \int_{-b}^{b} -f(u) \, du = -\int_{-b}^{b} f(u) \, du.
\]
Esto implica que $\displaystyle \int_{-b}^{b} f(x) \, dx = 0$.

%---
%\section*{Ejemplo con Función Par}
%Calcular $\int_{-2}^{2} x^2 \cos(x) \, dx$:
%\begin{itemize}
%    \item $f(x) = x^2 \cos(x)$ es \textbf{par} porque $f(-x) = (-x)^2 \cos(-x) = x^2 \cos(x) = f(x)$.
%    \item Por lo tanto:
%    \[
%    \int_{-2}^{2} x^2 \cos(x) \, dx = 2 \int_{0}^{2} x^2 \cos(x) \, dx.
%    \]
%\end{itemize}


\end{document}